\section{Related Work}\label{sec:related}

\paragraph{Payment channels and layer-2 protocols.}
The concept of off-chain payment channels dates to Spilman~\cite{spilman2013} and was substantially generalized by Poon and Dryja's Lightning Network~\cite{poon2016lightning}, which chains channels into a routed graph using HTLCs. Decker and Wattenhofer~\cite{decker2015fast} propose duplex micropayment channels using timelocked commitment trees. The \emph{eltoo} (ANYPREVOUT) update mechanism~\cite{decker2018eltoo} simplifies channel state management at the protocol level but still assumes a bidirectional two-party channel. Broader surveys~\cite{gudgeon2020sok} catalogue the full space of layer-two designs, and concurrent work has extended channel networks with privacy and concurrency guarantees~\cite{malavolta2017concurrency}, atomic multi-channel updates~\cite{egger2019atomic}, off-chain rebalancing~\cite{khalil2017revive}, and faster state-channel state machines~\cite{miller2019sprites}. All of these designs require the endpoints of a channel to be fixed at open time and to be online during payment. \thunderbolt differs by allowing the ``endpoint'' to change: the holder role moves from one user to the next via off-chain transfer, without opening a new channel.

\paragraph{Virtual channels and channel factories.}
Dziembowski et al.~\cite{dziembowski2018virtual} introduce virtual channels that avoid the need for a direct link between payer and payee by tunneling through an intermediary's channels. Burchert et al.~\cite{burchert2018scalable} propose channel factories that batch-open many channels in a single on-chain transaction. Aumayr et al.~\cite{aumayr2021bitcoin} adapt the virtual-channel idea to Bitcoin's limited scripting environment. These reduce per-channel on-chain cost but still require channels between specific pairs; \thunderbolt eliminates per-pair channels entirely via a single UTXO lock supporting unbounded holder transitions.

\paragraph{Verifiable encryption of discrete logarithms.}
Camenisch and Shoup~\cite{camenisch2003verifiable} give an efficient Sigma-protocol~\cite{cramer1997modular, damgard2002sigma} proof that a Paillier ciphertext encrypts the discrete log of a known group element, now standard in threshold cryptography, anonymous credentials~\cite{camenisch2000verifiable}, voting, and escrow. Asokan et al.~\cite{asokan1998optimistic} give earlier constructions for verifiable encryption in the context of optimistic fair exchange. Our Invoice and Transfer Consistency Proofs are direct instantiations of Camenisch--Shoup verifiable encryption, specialized to the ``one add, one subtract'' invariant of \thunderbolt.

\paragraph{Verifiable encrypted secret sharing.}
Shoup's Simple VESS~\cite{shoup2024simple} provides a verifiable encryption of Shamir secret sharing~\cite{shamir1979share} to a committee, designed as a building block for distributed key generation~\cite{gennaro2007secure} and resharing. Earlier work on verifiable secret sharing includes Feldman~\cite{feldman1987practical} and Pedersen~\cite{pedersen1991non}, which provide information-theoretic verification but not encryption to specific recipients. \thunderbolt uses VESS when the Committee is instantiated as a threshold group, to distribute the committee-side refresh $r_s$ so that no single member learns the committee share $y_s$.

\paragraph{Threshold signatures for Bitcoin.}
FROST~\cite{komlo2021frost} and ROAST~\cite{ruffing2022roast} provide efficient threshold Schnorr~\cite{schnorr1991efficient} signing, which can be used to finalize a \thunderbolt channel when the on-chain output is a Taproot~\cite{bip341} key. Lindell~\cite{lindell2017fast} gives a two-party ECDSA protocol, and Canetti et al.~\cite{canetti2020uc} a proactive threshold-ECDSA protocol with UC security. Compact multi-signatures~\cite{boneh2018compact} offer a complementary route to reducing on-chain key material. These schemes are complementary to \thunderbolt: they handle the on-chain finalization, while our Sigma-protocol proofs handle the off-chain state transitions.

\paragraph{General-purpose zero-knowledge proofs.}
Groth16~\cite{groth16}, Plonk~\cite{gabizon2019plonk}, transparent STARKs~\cite{ben2019scalable}, and zkVM-based systems such as Succinct SP1~\cite{sp1} and RISC Zero~\cite{risczero} can in principle prove any NP statement. While powerful, these systems pay a large constant for algebraic relations that Sigma protocols handle natively. Recursive proof composition without trusted setup~\cite{bowe2019halo} and accumulation-based aggregation~\cite{boneh2020recursive} reduce verification costs at scale but do not eliminate the proving-time overhead for Paillier-heavy statements.

\paragraph{Off-chain protocols beyond Bitcoin.}
Ethereum's layer-2 landscape includes optimistic rollups~\cite{arbitrum} and zk-rollups~\cite{zkSync}, which batch many transactions into a single on-chain proof or challenge. These systems benefit from Ethereum's expressive smart contracts and are not directly comparable to Bitcoin-native protocols. \thunderbolt is designed specifically for Bitcoin's constrained scripting environment.
